\documentclass[a4paper,10pt]{article}
\usepackage[utf8]{inputenc}
\usepackage[french] {babel}
\usepackage[T1]{fontenc}
\usepackage{lmodern}
\usepackage{graphicx}
\usepackage{graphics}
\usepackage{ulem}
\usepackage{amssymb}
\usepackage{url}
\usepackage[a4paper]{geometry}
\geometry{hscale=0.7,vscale=0.7,centering}
\usepackage{vmargin}
\usepackage{amsmath}
\usepackage{amssymb}
\usepackage{amsthm}
\usepackage{moreverb}
\usepackage{listings}
\newtheorem{theorem}{Théorème}[section]
\newtheorem{defi}{Définition}[section] 
\newtheorem{prop}{Propriété}[section] 
\usepackage{color}
\definecolor{gris}{rgb}{0.95,0.95,0.95}
\lstset{numbers=left, tabsize=4, backgroundcolor=\color{gris},
frame=single, breaklines=true,
keywordstyle=\color{black},
stringstyle=\ttfamily,
framexleftmargin=6mm, xleftmargin=6mm}
%opening
\title{LINGI 2261 : Artificial Intelligence \\
Assignement 3 - Mid-Project}
\author{Rochet Florentin - Debroux Léonard} 
\date{Année académique 2011-2012}

\begin{document}

	\begin{titlepage}
		\begin{center}
			{\huge LINGI2261: Artificial Intelligence}\\
			\vspace{0.4cm}
			
			{\Large {Professor : Yves Deville\\ \vspace{0.2cm} Teaching assistants : Cyrille Dejemeppe and Jean-Baptiste Mairy  }}\\
			\vspace{0.6cm}
			
			{\Large \textit{ Assignement3 : Adversarial Search}}\\
			\vspace{1.2cm}

			\texttt{}\\
			\vspace{0.2cm}

			\includegraphics[height=10cm]{pageGarde.png}\\
			\vspace{0.1cm}
			{\Large \textbf{Universit\'e Catholique de Louvain}}
			\vspace{0.7cm}

			Groupe 37 \\
			\vspace{0.2cm}
			
			Florentin Rochet \\
			Léonard Debroux\\
			\vspace{0.2cm}
			2012-2013\\
		\end{center}
	\end{titlepage}

	\newpage
	

	\section{Scipion}
	
	\section{Alpha-Beta search}	
	
	\section{Sarena}
		\subsection{A basic Alpha-Beta agent}
		
		\subsection{Comparison of MiniMax and Alpha-Beta}
	
	\section{Search Algorithms}
		\subsection{Statments to be proven}
			\subsubsection{Question 1: Prove that depth-first is a special case of uniform-cost search}
				We have to find an example of a tree that have the same search-path for the depth-first and the uniform-cost. If a tree is build is such as the cost of the nodes increments by one following the path of an actual depth-search, the uniform-cost will follow that same path and thus perform a depth-first search. 
				
			\subsubsection{Question 2: Prove that breadth-first is a special case of uniform-cost search}	
				In the simple case where all the nodes have the exact same cost, the uniform-search will run through the tree by level exactly like a breadth-first search.
			\subsubsection{Question 3: Prove that greedy-best-first search is a special case of A*}
				The f(n) function of an A* search is g(n) + h(n), which is the sum of the exact cost from the beginning to a node n and the extimated cost provided by the heuristic. In a greedy-best first, g(n) is equal to zero and thus it is a special case of A* where we do not take into account the cost to get to the node n. 
			\subsubsection{Question 4: Prove that uniform-cost is a special case of A*}
				This is pretty much the same idea than Q3, only here, it is h(n) that is equal to zero which means that we do not use a heuristic. f(n) is equal only to g(n) that represent the cost from the beginning to the node n.\\
			
			\subsubsection{Question 5: Prove that if a heuristic is consistent, then it is admissible}
			
				\begin{proof}
					Let's procede by in a recursive way, using the consitency definition :						$$ h(n) \leq c(n,a,n') + h(n') $$
					We assume first that n' is the goal, and thus, h(g) equals 0 by definition. We get:
					$$ h(n) \leq c(n,a,g) $$
					Where n is the node strictly before reaching the goal. We see that h is admissible for this node. We can now start the recursion like so:
					$$ h(n) \leq c(n,a,n') + h(n') $$
					Knowing that h(n') is admissible, we can set an upper bound for h(n') that is the exact minimal cost to the goal and can write the above as:
					$$ h(n) \leq c(n,a,n') + C*(n', g) $$
					The right term can thus be written as:
					$$ h(n) \leq C*(n, g) $$
					Which is the definition of admissibility.        
				\end{proof}
				We have thus prove that a consistent heuristic is admissible.
				
		\subsection{A* versus uniform-cost search}
			\subsubsection{Question 1: Give a consistent heuristic for this problem. Prove that it is admissible}
				A consistent heuristic for this problem is to use h(n) = manhattan distance from n to the goal. We will prove that it is consistent and thanks to the result of question 1.5, we are sure that it will be admissible too.
				\begin{proof}
					To prove that this heuristic is consistent and thus, admissible, we consider that the consistency statement is false and get:
					$$ h(n) > c(n,a,n') + h(n') $$
					Where "a" is a move in any of the cardinal directions.
					By subtracting h(n') on both side we get:
					$$ h(n) - h(n') > c(n,a,n') $$
					The manhattan distance heuristic should have a cost of 1 per action and accordingly to the latter expression, it is at least of 2 because c(n,a,n') is equal to 1.\\
					This is impossible and thus, the statement of consistency must be true.\\
					Because the heuristic is consistent, it is admissible as proved earlier.   
				\end{proof}  				
				
			\subsubsection{Question 2: Show on the left maze the states (board positions) that are visited during an execution of A* graph search with a manhattan distance heuristic (ignoring walls). A state is visited when it is selected in the fringe and expanded. In the A* algorithm, we assume that when different states in the fringe have the smallest value, the algorithm chooses the state with the smallest coordinate (i, j) ((0, 0) being the bottom left position, i being the horizontal index and j the vertical one) using a lexicographical order}
				For the following search, when looking at the children of a node, we start looking to the north and continue clockwise.\\
			\begin{center}
				\begin{tabular}{|c|c|c|c|c|c|c|}
					\hline 
					  &    &    &    &    &    &   \\ 
					\hline 
					  &    &    &    &    &    &   \\ 
					\hline 
					  &    &    &    &    &    &   \\ 
					\hline 
					  &    &    &    &    &    &   \\ 
					\hline 
					9 & 10 & 11 & 12 &    &    &   \\ 
					\hline 
					7 & 8  & \# & 13 &    &    &   \\ 
					\hline 
					5 & 6  & \# & 14 & 17 & 19 &   \\ 
					\hline 
					2 & 4  & \# & 15 & 16 & 18 & 20 \\ 
					\hline 
					1 & 3  & \# &    &    &    &   \\ 
					\hline 
				\end{tabular}
			\end{center} 
				The path that is retained by the A* search using the manhattan distance heuristic is the following: \{1, 2, 4, 6, 8, 10, 11, 12, 13, 14, 15, 16, 18, 20\}. The numbers indicates in which order the nodes are seen.

			\subsubsection{Show on the right maze the board positions visited by a uniform-cost graph search. When several states have the smallest path cost, this uniform-cost search visits them in the same lexicographical order as A*}
			 \begin{center}
				\begin{tabular}{|c|c|c|c|c|c|c|}
					\hline 
					19 & 25 & 32 & 41 & 49 & 55 &    \\ 
					\hline 
					15 & 20 & 26 & 24 & 43 & 51 &    \\ 
					\hline 
					12 & 16 & 21 & 28 & 36 & 45 & 53 \\ 
					\hline 
					10 & 13 & 17 & 23 & 30 & 38 & 47 \\ 
					\hline 
					8  & 11 & 14 & 18 & 24 & 31 & 39 \\ 
					\hline 
					6  & 9  & \# & 22 & 29 & 37 & 46 \\ 
					\hline 
					4  & 7  & \# & 27 & 35 & 44 & 52 \\ 
					\hline 
					2  & 5  & \# & 33 & 42 & 50 & 56 \\ 
					\hline 
					1  & 3  & \# & 40 & 48 & 54 &    \\ 
					\hline 
				\end{tabular} 
			\end{center}
				The uniform cost search is less effective than the A* using a manhattan distance, though the path is also of length 13. That path follows the following set of nodes: \{1, 2, 4, 6, 8, 11, 14, 18, 22, 27, 33, 42, 50, 56\}
			
	\newpage
	
	\section{Sokoban planning problem}
	
		\subsection{Question 1: As illustrated on Figure 3 some situations cannot lead to a solution. Are there other similar situations? If yes, describe them.}
			
			There is another situation that cannot lead to a solution, that is when a block is against a wall and must go in a perpendicular direction w.r.t. that wall as shown on the representation below: \\

	\begin{center}
	\begin{tabular}{ccccc}
		\# & \# & \# & \# & \# \\ 
		\# &    &    &    & \# \\ 
		\# & .  &    & \$ & \# \\ 
		\# &    &    &    & \# \\ 
		\# & \# & \# & \# & \# \\ 
	\end{tabular}
	\end{center}
	
	In this case, the box is stuck against the right wall and it will never be possible to push it to the left to reach the goal.\\
	In the case of a goal being against a wall, there is a dead state if two boxes a pushed against that wall as shown on the example below: \\
	
	\begin{center}
	\begin{tabular}{ccccccccc}
		\# & \# & \# & \# & \# & \# & \# & \# & \# \\ 
		\# &    & .  &    & \$ &    & \$ &    & \# \\ 
		\# &    &    &    &    &    &    &    & \# \\ 
		\# &    &    &    &    &    &    &    & \# \\ 
		\# & \# & \# & \# & \# & \# & \# & \# & \# \\ 
	\end{tabular}
	\end{center}
	The box on the right will be able to go on the goal but the one on the right is in the same position as explain above, although there is a goal against the same wall as it.
		
		\subsection{Question 2: Why is it important to identify dead states in your successor function? How are you going to implement it?}
			
			It is important to identify the dead state because it avoids to generate and explore lots of useless states.\\
			 Let's assume that one box is on a dead state case, and two boxes can still be moved. Let's also say that there is about 20 reachable cases for the boxes left and for the character, the number of useless states would be around $ 20^{3} $ (not exactly) for only one box on one dead state. This is what shows the following example.\\

	\begin{center}
	\begin{tabular}{ccccccc}
		\# & \# & \# & \# & \# & \# & \# \\ 
		\# & \$ &    &    &    &    & \# \\ 
		\# &    & @  &    & \$ &    & \# \\ 
		\# &    &    & \$ &    &    & \# \\ 
		\# &    &    &    &    &    & \# \\
		\# & \# & \# & \# & \# & \# & \# \\ 
	\end{tabular}\\ 
	\end{center}

	
		The upper left box can not be moved anymore.\\
			To avoid that kind of state, we created two type of dead state. Both are detected at the beginning of the execution but do not imply the same behaviour.
			
			\subsubsection{Static Dead States}
				If the avatar is trying to move a box on a static dead state, it is directly forbidden and the move is not allowed, it is especially the case in corner, as show on the example right above.
				Static dead cases exist also in other case but this will be covered in the next section.
				
			\subsubsection{Possible Dead States}
				Those are created when there is a goal next to a wall (see example 2 in 2.1). When there is no goal along a wall, then all the cases become static dead states, when there is one or several, they become possible dead states. This means that when there is as much boxes as goals along that wall there can not be one more. It avoids similar scenario's to example 2 in 2.1.\\
				It is implemented thanks to a list currentDeadStates proper to every state, it contains the coordinates of every possible dead state that has turned to dead state because of the boxes on it.\\
				The move actions update this list and allow moves for boxes that are on it.				
				
		\subsection{Question 3: Describe possible (non trivial) heuristic(s) to reach a goal state (with reference if any). Is(are) your heuristic(s) admissible and/or consistent?}

		Except the trivial heuristic, we developed three heuristics, the one used in the question 5 for tests is described later in the document.\\
		
		The second non trivial is the following:\\ 
		The h(n) function minimizes the sum of the manhattan distance for the set of possibles configurations from boxes to goals (one box associated to one goal and one goal is targeted by one and only one box). \\
		So, if we have N boxes/goals on our problem, we have factorial N possible pair of box and goal. We take the minimum of the sum to ensure that, first the heuristic will never overestimate the cost (since we have the pure minimal cost to move the boxes to the goals). Then the heuristic is consistent because for a state n' closer to the goal from a previous state n, we have $ h(n) \geq h(n')$ by our definition of $h$ and more precisely the difference between the two if any is of $1$, we have then $h(n) \leq c(n,a,n')+h(n')$ with c the cost from $n$ to $n'$ by the action $a$ (which is an action towards the goal and not in the other way). So our heuristic is admissible and consistent.
		
The third that we developed count the amount of boxes not on a goal. It is not very accurate, but it provide an indication concerning the advancement of towards the solution. It is obviously admissible and consistent, the thought behind the prove is the same as seen in the course for the 8-puzzle.
		
		\subsection{Question 4: Implement this problem. Extend the Problem class and implement the necessary methods and other class(es) if necessary. Your file must be named sokoban.py. You program must print to the standard output a solution to the sokoban instance given in argument satisfying the described format.}
			The resolution is launched using python3 sokoban.py sokoInstxy with $x = \{0,1\}$ and $y = \{0 \ to \ 9\}$\\
			There is no need to add the extensions.

		\subsection{Question 5: Experiment, compare and analyze informed (astar\_graph \_search) and uninformed (breadth\_first\_graph\_search) graph search of aima-python3 on the 15 instances of sokoban provided. Report in a table the time, the number of explored nodes and the number of steps to reach the solution. Are the number of explored nodes always smaller with astar\_graph\_search, why? When no solution can be found by a strategy in a reasonable time (say 5 min), explain the reason (time-out and/or swap of the memory).}
			We implemented several heuristics and will use one of them for the tests. \\
			The one we will use the sum of the manhattan distance from each box to its closest goal (which means that a goals can be wanted by several boxes). \\
			This heuristic is consistent because it is the sum on the box to their closest goals. As seen in the course, the manhattan distance from a tile coordinate to its goal is consistent. Because of the fact that we have as much terms in our sum as we have boxes, it is the minimum amount of moves that are to be done to reach a solution in the best case in which several boxes could go on the same goal. It is obviously underestimated the exact cost and making. And we respect the consistency because of the use of the manhattan distance as proved in section 1.2.1.\\
			All the tests were run on the same computer to ensure consistency of the results.\\
			
			\begin{tabular}{|c||c||c|c||c|c|}
				\hline 
				\textbf{Problem} & \textbf{Steps} & \textbf{BFS Time [s]} & \textbf{BFS Nodes} & \textbf{A* Time [s]} & \textbf{A* Nodes} \\ 
				\hline 
				sokoInst01 & 15 & 0.045 & 1,548 & 0.029 & 707 \\ 
				\hline
				sokoInst02 & 65 & 0.152 & 6,922 & 0.215 & 6,859 \\ 
				\hline
				sokoInst03 & 135 & 6.963 & 238,760 & 9.553 & 237,686 \\ 
				\hline
				sokoInst04 & 26 & 0.736 & 27,741 & 1.003 & 21,305 \\ 
				\hline
				sokoInst05 & 108 & 0.435 & 19,729 & 0.637 & 19,724 \\ 
				\hline
				sokoInst06 & 27 & 1.624 & 52,405 & 2.85 & 45,774 \\ 
				\hline
				sokoInst07 & 98 & 0.092 & 4,770 & 0.118 & 4,439 \\ 
				\hline
				sokoInst08 & 90 & 0.138 & 6,323 & 0.195 & 6,290 \\ 
				\hline
				sokoInst09 & 101 & 2.333 & 84,559 & 3.127 & 83,401 \\ 
				\hline
				sokoInst10 & 42 & 0.079 & 3,268 & 0.099 & 2,832 \\ 
				\hline
				sokoInst11 & 179 & 7.5 & 237,430 & 10.079 & 231,116 \\ 
				\hline
				sokoInst12 & 70 & 3.819 & 129,777 & 4.375 & 104,220 \\ 
				\hline
				sokoInst13 & 182 & 3.26 & 127,068 & 4.547 & 126,693 \\ 
				\hline
				sokoInst14 & 65 & 4.839 & 157,453 & 7.042 & 150,731 \\ 
				\hline
				sokoInst15 & 56 & 3.16 & 90,014 & 3.668 & 78,777 \\ 
				\hline
			\end{tabular}\\
			
			We can see that using A* with a heuristic lead to explore less nodes, but we also see a slight increase in the time of computation. This comes from the fact that avoiding to explore a node cost more time due to the computation of the heuristic.\\
			The number of nodes is always smaller using A* search because in the worse case, A* would perform a BFS, so, in any other case, there is less nodes explored.
			
		\subsection{Question 6: What are the performances of your program when you don’t perform dead state detection?}
			The answer to this question is theoretically that the number of nodes would explode and so that the computation time would also increase. To have a better idea of those values, we ran the 15 instances of sokoban provided without detecting the dead states. //
			
			The time taken is a bit higher than should be, because we suppressed the dead states detection, but the search still goes through some functions especially designed for dead states management become useless. The process to get rid of those is heavy and the relevant information here is the amount of nodes explored which is not affected by this choice.\\
			The following array is similar to the previous one but without dead states management:\\
			
			\begin{tabular}{|c||c||c|c||c|c|}
				\hline 
				\textbf{Problem} & \textbf{Steps} & \textbf{BFS Time [s]} & \textbf{BFS Nodes} & \textbf{A* Time [s]} & \textbf{A* Nodes} \\ 
				\hline 
				sokoInst01 & 15 & 0.06 & 2,409 & 0.033 & 833 \\ 
				\hline
				sokoInst02 & 65 & 0.68 & 30,453 & 0.99 & 29,609 \\ 
				\hline
				sokoInst03 & 135 & 56.07 & 1,527,231 & 93.83 & 1,524,338 \\ 
				\hline
				sokoInst04 & 26 & 43.18 & 1,119,877 & 82.09 & 730,806 \\ 
				\hline
				sokoInst05 & 108 & 4.224 & 139,784 & 5.919 & 139,603 \\ 
				\hline
				sokoInst06 & 27 & None & None & 1962.7 & 2,868,224 \\ 
				\hline
				sokoInst07 & 98 & 0.212 & 9,392 & 0.276 & 8,993 \\ 
				\hline
				sokoInst08 & 90 & 0.293 & 13,626 & 0.486 & 13,348 \\ 
				\hline
				sokoInst09 & 101 & 7.348 & 241,258 & 11.06 & 239,760 \\ 
				\hline
				sokoInst10 & 42 & 0.632 & 24,058 & 0.701 & 19,038 \\ 
				\hline
				sokoInst11 & 179 & 25.35 & 713,439 & 35.62 & 686,897 \\ 
				\hline
				sokoInst12 & 70 & 10.12 & 287,851 & 11.14 & 237,661 \\ 
				\hline
				sokoInst13 & 182 & 15.08 & 419,669 & 18.38 & 417,814 \\ 
				\hline
				sokoInst14 & 65 & 47.27 & 1,332,657 & 99.08 & 1,306,553 \\ 
				\hline
				sokoInst15 & 56 & 29.29 & 802,991 & 50.8 & 739,200 \\ 
				\hline
			\end{tabular}\\
			
			The "None" statement means that the search filled the memory so that the OS killed the process, this is due to the huge amount of states generated before reaching a solution. The reason for it is that all the cases along the wall in the 6th configuration are in dead states when detecting them, so instead of 30 cases where boxes can go, there is only 12. The exponential growth for that kind of problem explains that the search does not provide a result.\\
			When using a heuristic, the memory is filled at a slower rate which allow the search to find a solution. The result of about 32 minutes is given as an indication.\\
			The reason it takes so much time is the same as the reason the BFS doesn't find a solution.
			
			We can see by analysing that array that the number of explored nodes explodes and so does the time (but as said before, it is not the most relevant information).
			
\end{document}